\documentclass{article}
\usepackage[T1]{fontenc}
\usepackage{lmodern}
\usepackage{graphicx}
\usepackage{amsmath,amssymb}
\usepackage[a4paper,margin=2cm]{geometry}
\usepackage[absolute,overlay]{textpos}
\setlength{\TPHorizModule}{1mm}
\setlength{\TPVertModule}{1mm}
\usepackage[indonesian]{babel}
\usepackage{hyperref}
\setlength{\emergencystretch}{2em}
% unit: Halaman 51

\begin{document}

% === Halaman 51 (python/native) ===

51

• Secara intuitif: jika jarak antara dua buah string yang berulang di 

dalam plainteks merupakan kelipatan dari panjang kunci, 

• maka string yang sama tersebut akan muncul menjadi kriptogram 

yang sama pula di dalam cipherteks. 

• Pada Contoh 1, 


- kunci = abcd 


- panjang kunci = 4


- jarak antara dua crypto yang berulang = 16


- 16 = kelipatan 4

  crypto dienkripsi menjadi kriptogram yang sama

Plainteks   : cryptoisshortforcryptography

Kunci 

   : abcdabcdabcdabcdabcdabcdabcd

Cipherteks: CSASTPKVSIQUTGQUCSASTPIUAQJB 

16

\end{document}
